\documentclass[tikz,border=4pt,12pt]{standalone}
\usepackage{ifthen}
\usetikzlibrary{arrows.meta,calc,positioning}
% Shared TikZ macros for the architecture figures (Medverse vs. Ours).
% Load with % Shared TikZ macros for the architecture figures (Medverse vs. Ours).
% Load with % Shared TikZ macros for the architecture figures (Medverse vs. Ours).
% Load with \input{arch-macros.tex} after \usepackage{tikz,ifthen} and
% \usetikzlibrary{arrows.meta,calc,positioning}. All macro names are
% prefixed `arch` to avoid clashing with pgf/tikz internals.

% ---- palette (approximating the PowerPoint drafts) ----
\definecolor{archTargetGreen}{HTML}{9BC7AC}
\definecolor{archContextPink}{HTML}{E8B8B0}
\definecolor{archAttnBlue}{HTML}{BBD6EC}
\definecolor{archEncGray}{HTML}{E4E4E4}
\definecolor{archArrowPale}{HTML}{CFE6D6}
\definecolor{archCtBg}{HTML}{2B2B2B}
\definecolor{archCtBlob}{HTML}{8C8C8C}
\definecolor{archCtxYellow}{HTML}{E8D24C}
% darker, more saturated siblings of archTargetGreen/archContextPink --
% pastel fills read poorly as 1pt strokes, and mask/label content needs to
% read as "which volume owns this" at a glance, so masks/labels and any
% line (arrow, brace) touching that volume's row use these instead of the
% pastel role color or the old flat archMaskGold.
\definecolor{archTargetAccent}{HTML}{2F6B4A}  % target: darker green
\definecolor{archContextAccent}{HTML}{A8453E} % context: darker red

% fake-isometric depth skew shared by every cube
\newcommand{\archdx}{0.30}
\newcommand{\archdy}{0.20}

% ---- shared text-role macros ----
% Every arch_*.tex figure is compiled standalone at the SAME base size
% (documentclass[tikz,border=4pt,12pt]{standalone}, matching the thesis
% body's 12pt), so \normalsize etc. mean the same absolute point size in
% every figure. Using these three named roles instead of ad hoc \Large /
% \bfseries\normalsize / \scriptsize per figure keeps "a box title" (etc.)
% the same size everywhere it appears, regardless of which figure it's in
% or how that figure's canvas happens to be scaled by \includegraphics.
%   \archtitle    -- figure/method/panel titles, the "xL" repeat marker
%   \archboxlabel -- sub-block header labels (Fuse, Encoder, attention, ...)
%   \archnote     -- small secondary annotations (placeholder captions)
% Plain, unlabeled node text (cube/token labels, row/lane labels) stays at
% the bare \normalsize default -- no macro needed for that role.
%
% DECLARATIONS, not argument-taking macros -- deliberately, so they can
% precede a raw `\\` inside a tikz node's text (as `{\archboxlabel line
% one\\line two}`) without hiding that `\\` one macro-expansion level deep.
% A `\\` buried inside a wrapped {\archboxlabel{...}} argument breaks
% TikZ's own align=center line-splitting (it scans for `\\` before the
% argument expands) -- use these the same way `\bfseries`/`\Large` are
% already used directly inline elsewhere in these figures.
\newcommand{\archtitle}{\Large\bfseries}
\newcommand{\archboxlabel}{\bfseries\normalsize}
\newcommand{\archnote}{\scriptsize}

% ---------------------------------------------------------------
% \archvolcube{x}{y}{w}{h}{topcolor}{content}{label}
%   content: ct | ctmask | predmask | <flat color name>
% ---------------------------------------------------------------
\newcommand{\archvolcube}[7]{%
  \begin{scope}
    \pgfmathsetmacro{\cx}{#1}\pgfmathsetmacro{\cy}{#2}%
    \pgfmathsetmacro{\w}{#3}\pgfmathsetmacro{\h}{#4}%
    \fill[#5] (\cx,\cy+\h) -- ++(\archdx,\archdy) -- ++(\w,0) -- ++(-\archdx,-\archdy) -- cycle;
    \fill[#5!60!black] (\cx+\w,\cy) -- ++(\archdx,\archdy) -- ++(0,\h) -- ++(-\archdx,-\archdy) -- cycle;
    \begin{scope}
      \clip (\cx,\cy) rectangle ++(\w,\h);
      \ifthenelse{\equal{#6}{ct}}{%
        \fill[archCtBg] (\cx,\cy) rectangle ++(\w,\h);
        \fill[archCtBlob] plot[smooth cycle] coordinates
          {(\cx+.25*\w,\cy+.55*\h) (\cx+.45*\w,\cy+.75*\h) (\cx+.7*\w,\cy+.6*\h)
           (\cx+.6*\w,\cy+.3*\h) (\cx+.35*\w,\cy+.25*\h)};
      }{\ifthenelse{\equal{#6}{ctmask}}{%
        \fill[archCtBg] (\cx,\cy) rectangle ++(\w,\h);
        \fill[archCtBlob] plot[smooth cycle] coordinates
          {(\cx+.25*\w,\cy+.55*\h) (\cx+.45*\w,\cy+.75*\h) (\cx+.7*\w,\cy+.6*\h)
           (\cx+.6*\w,\cy+.3*\h) (\cx+.35*\w,\cy+.25*\h)};
        \fill[archCtxYellow] plot[smooth cycle] coordinates
          {(\cx+.42*\w,\cy+.42*\h) (\cx+.52*\w,\cy+.5*\h) (\cx+.5*\w,\cy+.35*\h) (\cx+.44*\w,\cy+.3*\h)};
      }{\ifthenelse{\equal{#6}{predmask}}{%
        \fill[black] (\cx,\cy) rectangle ++(\w,\h);
        \fill[white] plot[smooth cycle] coordinates
          {(\cx+.4*\w,\cy+.65*\h) (\cx+.6*\w,\cy+.7*\h) (\cx+.62*\w,\cy+.45*\h)
           (\cx+.48*\w,\cy+.3*\h) (\cx+.38*\w,\cy+.4*\h)};
      }{%
        \fill[#6] (\cx,\cy) rectangle ++(\w,\h);
      }}}%
    \end{scope}
    \draw[black,thin] (\cx,\cy) rectangle ++(\w,\h);
    \draw[black,thin] (\cx,\cy+\h) -- ++(\archdx,\archdy) -- ++(\w,0) -- ++(-\archdx,-\archdy) -- cycle;
    \draw[black,thin] (\cx+\w,\cy) -- ++(\archdx,\archdy) -- ++(0,\h) -- ++(-\archdx,-\archdy) -- cycle;
    \node[anchor=north,align=center] at (\cx+\w/2,\cy-0.1) {#7};
  \end{scope}
}

% ---------------------------------------------------------------
% \archfeatcube{x}{y}{w}{h}{color}{label} -- flat-colored feature volume
% ---------------------------------------------------------------
\newcommand{\archfeatcube}[6]{%
  \begin{scope}
    \pgfmathsetmacro{\cx}{#1}\pgfmathsetmacro{\cy}{#2}%
    \pgfmathsetmacro{\w}{#3}\pgfmathsetmacro{\h}{#4}%
    \fill[#5] (\cx,\cy) rectangle ++(\w,\h);
    \fill[#5!80!white] (\cx,\cy+\h) -- ++(\archdx,\archdy) -- ++(\w,0) -- ++(-\archdx,-\archdy) -- cycle;
    \fill[#5!60!black] (\cx+\w,\cy) -- ++(\archdx,\archdy) -- ++(0,\h) -- ++(-\archdx,-\archdy) -- cycle;
    \draw[black,thin] (\cx,\cy) rectangle ++(\w,\h);
    \draw[black,thin] (\cx,\cy+\h) -- ++(\archdx,\archdy) -- ++(\w,0) -- ++(-\archdx,-\archdy) -- cycle;
    \draw[black,thin] (\cx+\w,\cy) -- ++(\archdx,\archdy) -- ++(0,\h) -- ++(-\archdx,-\archdy) -- cycle;
    \node[align=center] at (\cx+\w/2,\cy+\h/2) {#6};
  \end{scope}
}

% ---------------------------------------------------------------
% \archfunnel{x}{y}{h}{wwide}{wnarrow}{dir}{color}{label}
%   dir=1  : wide on left, narrow on right  (encoder)
%   dir=-1 : narrow on left, wide on right  (decoder)
% ---------------------------------------------------------------
\newcommand{\archfunnel}[8]{%
  \begin{scope}
    \pgfmathsetmacro{\cx}{#1}\pgfmathsetmacro{\cy}{#2}\pgfmathsetmacro{\h}{#3}%
    \pgfmathsetmacro{\ww}{#4}\pgfmathsetmacro{\wn}{#5}%
    \pgfmathsetmacro{\inset}{(\ww-\wn)/2}%
    \ifnum#6=1
      \fill[#7] (\cx,\cy) -- (\cx+\ww,\cy+\inset) -- (\cx+\ww,\cy+\h-\inset) -- (\cx,\cy+\h) -- cycle;
      \draw[black,thin] (\cx,\cy) -- (\cx+\ww,\cy+\inset) -- (\cx+\ww,\cy+\h-\inset) -- (\cx,\cy+\h) -- cycle;
    \else
      \fill[#7] (\cx,\cy+\inset) -- (\cx+\ww,\cy) -- (\cx+\ww,\cy+\h) -- (\cx,\cy+\h-\inset) -- cycle;
      \draw[black,thin] (\cx,\cy+\inset) -- (\cx+\ww,\cy) -- (\cx+\ww,\cy+\h) -- (\cx,\cy+\h-\inset) -- cycle;
    \fi
    \node[align=center] at (\cx+\ww/2,\cy+\h/2) {#8};
  \end{scope}
}

% ---------------------------------------------------------------
% \archbowtie{x}{y}{w}{h}{pinch}{color}{label} -- Medverse-style
% single-row "dual U-Net" hourglass (funnel-in then funnel-out).
% ---------------------------------------------------------------
\newcommand{\archbowtie}[7]{%
  \begin{scope}
    \pgfmathsetmacro{\cx}{#1}\pgfmathsetmacro{\cy}{#2}%
    \pgfmathsetmacro{\w}{#3}\pgfmathsetmacro{\h}{#4}\pgfmathsetmacro{\p}{#5}%
    \pgfmathsetmacro{\inset}{(\h-\p)/2}%
    \fill[#6] (\cx,\cy) -- (\cx+\w/2,\cy+\inset) -- (\cx+\w,\cy) -- (\cx+\w,\cy+\h)
      -- (\cx+\w/2,\cy+\h-\inset) -- (\cx,\cy+\h) -- cycle;
    \draw[black,thin] (\cx,\cy) -- (\cx+\w/2,\cy+\inset) -- (\cx+\w,\cy) -- (\cx+\w,\cy+\h)
      -- (\cx+\w/2,\cy+\h-\inset) -- (\cx,\cy+\h) -- cycle;
    \node[align=center] at (\cx+\w/2,\cy+\h/2) {#7};
  \end{scope}
}

% ---------------------------------------------------------------
% \archtokengrid{x}{y}{rows}{cols}{cell}{color} -- grid of patch tokens
% ---------------------------------------------------------------
\newcommand{\archtokengrid}[6]{%
  \begin{scope}
    \pgfmathsetmacro{\cx}{#1}\pgfmathsetmacro{\cy}{#2}\pgfmathsetmacro{\cell}{#5}%
    \pgfmathtruncatemacro{\rmax}{#3-1}%
    \pgfmathtruncatemacro{\cmax}{#4-1}%
    \foreach \r in {0,...,\rmax}{%
      \foreach \c in {0,...,\cmax}{%
        \fill[#6] ({\cx+\c*\cell},{\cy+\r*\cell}) rectangle ++(\cell*0.82,\cell*0.82);
        \draw[black,thin] ({\cx+\c*\cell},{\cy+\r*\cell}) rectangle ++(\cell*0.82,\cell*0.82);
      }%
    }%
  \end{scope}
}

% ---------------------------------------------------------------
% \archarrow{start}{end}{color} -- block arrow between stages
% ---------------------------------------------------------------
% #1 and #2 are full coordinates including parentheses, e.g. (1,2)
\newcommand{\archarrow}[3]{%
  \draw[-{Latex[length=2.4mm,width=2mm]},line width=2.2pt,#3] #1 -- #2;
}
 after \usepackage{tikz,ifthen} and
% \usetikzlibrary{arrows.meta,calc,positioning}. All macro names are
% prefixed `arch` to avoid clashing with pgf/tikz internals.

% ---- palette (approximating the PowerPoint drafts) ----
\definecolor{archTargetGreen}{HTML}{9BC7AC}
\definecolor{archContextPink}{HTML}{E8B8B0}
\definecolor{archAttnBlue}{HTML}{BBD6EC}
\definecolor{archEncGray}{HTML}{E4E4E4}
\definecolor{archArrowPale}{HTML}{CFE6D6}
\definecolor{archCtBg}{HTML}{2B2B2B}
\definecolor{archCtBlob}{HTML}{8C8C8C}
\definecolor{archCtxYellow}{HTML}{E8D24C}
% darker, more saturated siblings of archTargetGreen/archContextPink --
% pastel fills read poorly as 1pt strokes, and mask/label content needs to
% read as "which volume owns this" at a glance, so masks/labels and any
% line (arrow, brace) touching that volume's row use these instead of the
% pastel role color or the old flat archMaskGold.
\definecolor{archTargetAccent}{HTML}{2F6B4A}  % target: darker green
\definecolor{archContextAccent}{HTML}{A8453E} % context: darker red

% fake-isometric depth skew shared by every cube
\newcommand{\archdx}{0.30}
\newcommand{\archdy}{0.20}

% ---- shared text-role macros ----
% Every arch_*.tex figure is compiled standalone at the SAME base size
% (documentclass[tikz,border=4pt,12pt]{standalone}, matching the thesis
% body's 12pt), so \normalsize etc. mean the same absolute point size in
% every figure. Using these three named roles instead of ad hoc \Large /
% \bfseries\normalsize / \scriptsize per figure keeps "a box title" (etc.)
% the same size everywhere it appears, regardless of which figure it's in
% or how that figure's canvas happens to be scaled by \includegraphics.
%   \archtitle    -- figure/method/panel titles, the "xL" repeat marker
%   \archboxlabel -- sub-block header labels (Fuse, Encoder, attention, ...)
%   \archnote     -- small secondary annotations (placeholder captions)
% Plain, unlabeled node text (cube/token labels, row/lane labels) stays at
% the bare \normalsize default -- no macro needed for that role.
%
% DECLARATIONS, not argument-taking macros -- deliberately, so they can
% precede a raw `\\` inside a tikz node's text (as `{\archboxlabel line
% one\\line two}`) without hiding that `\\` one macro-expansion level deep.
% A `\\` buried inside a wrapped {\archboxlabel{...}} argument breaks
% TikZ's own align=center line-splitting (it scans for `\\` before the
% argument expands) -- use these the same way `\bfseries`/`\Large` are
% already used directly inline elsewhere in these figures.
\newcommand{\archtitle}{\Large\bfseries}
\newcommand{\archboxlabel}{\bfseries\normalsize}
\newcommand{\archnote}{\scriptsize}

% ---------------------------------------------------------------
% \archvolcube{x}{y}{w}{h}{topcolor}{content}{label}
%   content: ct | ctmask | predmask | <flat color name>
% ---------------------------------------------------------------
\newcommand{\archvolcube}[7]{%
  \begin{scope}
    \pgfmathsetmacro{\cx}{#1}\pgfmathsetmacro{\cy}{#2}%
    \pgfmathsetmacro{\w}{#3}\pgfmathsetmacro{\h}{#4}%
    \fill[#5] (\cx,\cy+\h) -- ++(\archdx,\archdy) -- ++(\w,0) -- ++(-\archdx,-\archdy) -- cycle;
    \fill[#5!60!black] (\cx+\w,\cy) -- ++(\archdx,\archdy) -- ++(0,\h) -- ++(-\archdx,-\archdy) -- cycle;
    \begin{scope}
      \clip (\cx,\cy) rectangle ++(\w,\h);
      \ifthenelse{\equal{#6}{ct}}{%
        \fill[archCtBg] (\cx,\cy) rectangle ++(\w,\h);
        \fill[archCtBlob] plot[smooth cycle] coordinates
          {(\cx+.25*\w,\cy+.55*\h) (\cx+.45*\w,\cy+.75*\h) (\cx+.7*\w,\cy+.6*\h)
           (\cx+.6*\w,\cy+.3*\h) (\cx+.35*\w,\cy+.25*\h)};
      }{\ifthenelse{\equal{#6}{ctmask}}{%
        \fill[archCtBg] (\cx,\cy) rectangle ++(\w,\h);
        \fill[archCtBlob] plot[smooth cycle] coordinates
          {(\cx+.25*\w,\cy+.55*\h) (\cx+.45*\w,\cy+.75*\h) (\cx+.7*\w,\cy+.6*\h)
           (\cx+.6*\w,\cy+.3*\h) (\cx+.35*\w,\cy+.25*\h)};
        \fill[archCtxYellow] plot[smooth cycle] coordinates
          {(\cx+.42*\w,\cy+.42*\h) (\cx+.52*\w,\cy+.5*\h) (\cx+.5*\w,\cy+.35*\h) (\cx+.44*\w,\cy+.3*\h)};
      }{\ifthenelse{\equal{#6}{predmask}}{%
        \fill[black] (\cx,\cy) rectangle ++(\w,\h);
        \fill[white] plot[smooth cycle] coordinates
          {(\cx+.4*\w,\cy+.65*\h) (\cx+.6*\w,\cy+.7*\h) (\cx+.62*\w,\cy+.45*\h)
           (\cx+.48*\w,\cy+.3*\h) (\cx+.38*\w,\cy+.4*\h)};
      }{%
        \fill[#6] (\cx,\cy) rectangle ++(\w,\h);
      }}}%
    \end{scope}
    \draw[black,thin] (\cx,\cy) rectangle ++(\w,\h);
    \draw[black,thin] (\cx,\cy+\h) -- ++(\archdx,\archdy) -- ++(\w,0) -- ++(-\archdx,-\archdy) -- cycle;
    \draw[black,thin] (\cx+\w,\cy) -- ++(\archdx,\archdy) -- ++(0,\h) -- ++(-\archdx,-\archdy) -- cycle;
    \node[anchor=north,align=center] at (\cx+\w/2,\cy-0.1) {#7};
  \end{scope}
}

% ---------------------------------------------------------------
% \archfeatcube{x}{y}{w}{h}{color}{label} -- flat-colored feature volume
% ---------------------------------------------------------------
\newcommand{\archfeatcube}[6]{%
  \begin{scope}
    \pgfmathsetmacro{\cx}{#1}\pgfmathsetmacro{\cy}{#2}%
    \pgfmathsetmacro{\w}{#3}\pgfmathsetmacro{\h}{#4}%
    \fill[#5] (\cx,\cy) rectangle ++(\w,\h);
    \fill[#5!80!white] (\cx,\cy+\h) -- ++(\archdx,\archdy) -- ++(\w,0) -- ++(-\archdx,-\archdy) -- cycle;
    \fill[#5!60!black] (\cx+\w,\cy) -- ++(\archdx,\archdy) -- ++(0,\h) -- ++(-\archdx,-\archdy) -- cycle;
    \draw[black,thin] (\cx,\cy) rectangle ++(\w,\h);
    \draw[black,thin] (\cx,\cy+\h) -- ++(\archdx,\archdy) -- ++(\w,0) -- ++(-\archdx,-\archdy) -- cycle;
    \draw[black,thin] (\cx+\w,\cy) -- ++(\archdx,\archdy) -- ++(0,\h) -- ++(-\archdx,-\archdy) -- cycle;
    \node[align=center] at (\cx+\w/2,\cy+\h/2) {#6};
  \end{scope}
}

% ---------------------------------------------------------------
% \archfunnel{x}{y}{h}{wwide}{wnarrow}{dir}{color}{label}
%   dir=1  : wide on left, narrow on right  (encoder)
%   dir=-1 : narrow on left, wide on right  (decoder)
% ---------------------------------------------------------------
\newcommand{\archfunnel}[8]{%
  \begin{scope}
    \pgfmathsetmacro{\cx}{#1}\pgfmathsetmacro{\cy}{#2}\pgfmathsetmacro{\h}{#3}%
    \pgfmathsetmacro{\ww}{#4}\pgfmathsetmacro{\wn}{#5}%
    \pgfmathsetmacro{\inset}{(\ww-\wn)/2}%
    \ifnum#6=1
      \fill[#7] (\cx,\cy) -- (\cx+\ww,\cy+\inset) -- (\cx+\ww,\cy+\h-\inset) -- (\cx,\cy+\h) -- cycle;
      \draw[black,thin] (\cx,\cy) -- (\cx+\ww,\cy+\inset) -- (\cx+\ww,\cy+\h-\inset) -- (\cx,\cy+\h) -- cycle;
    \else
      \fill[#7] (\cx,\cy+\inset) -- (\cx+\ww,\cy) -- (\cx+\ww,\cy+\h) -- (\cx,\cy+\h-\inset) -- cycle;
      \draw[black,thin] (\cx,\cy+\inset) -- (\cx+\ww,\cy) -- (\cx+\ww,\cy+\h) -- (\cx,\cy+\h-\inset) -- cycle;
    \fi
    \node[align=center] at (\cx+\ww/2,\cy+\h/2) {#8};
  \end{scope}
}

% ---------------------------------------------------------------
% \archbowtie{x}{y}{w}{h}{pinch}{color}{label} -- Medverse-style
% single-row "dual U-Net" hourglass (funnel-in then funnel-out).
% ---------------------------------------------------------------
\newcommand{\archbowtie}[7]{%
  \begin{scope}
    \pgfmathsetmacro{\cx}{#1}\pgfmathsetmacro{\cy}{#2}%
    \pgfmathsetmacro{\w}{#3}\pgfmathsetmacro{\h}{#4}\pgfmathsetmacro{\p}{#5}%
    \pgfmathsetmacro{\inset}{(\h-\p)/2}%
    \fill[#6] (\cx,\cy) -- (\cx+\w/2,\cy+\inset) -- (\cx+\w,\cy) -- (\cx+\w,\cy+\h)
      -- (\cx+\w/2,\cy+\h-\inset) -- (\cx,\cy+\h) -- cycle;
    \draw[black,thin] (\cx,\cy) -- (\cx+\w/2,\cy+\inset) -- (\cx+\w,\cy) -- (\cx+\w,\cy+\h)
      -- (\cx+\w/2,\cy+\h-\inset) -- (\cx,\cy+\h) -- cycle;
    \node[align=center] at (\cx+\w/2,\cy+\h/2) {#7};
  \end{scope}
}

% ---------------------------------------------------------------
% \archtokengrid{x}{y}{rows}{cols}{cell}{color} -- grid of patch tokens
% ---------------------------------------------------------------
\newcommand{\archtokengrid}[6]{%
  \begin{scope}
    \pgfmathsetmacro{\cx}{#1}\pgfmathsetmacro{\cy}{#2}\pgfmathsetmacro{\cell}{#5}%
    \pgfmathtruncatemacro{\rmax}{#3-1}%
    \pgfmathtruncatemacro{\cmax}{#4-1}%
    \foreach \r in {0,...,\rmax}{%
      \foreach \c in {0,...,\cmax}{%
        \fill[#6] ({\cx+\c*\cell},{\cy+\r*\cell}) rectangle ++(\cell*0.82,\cell*0.82);
        \draw[black,thin] ({\cx+\c*\cell},{\cy+\r*\cell}) rectangle ++(\cell*0.82,\cell*0.82);
      }%
    }%
  \end{scope}
}

% ---------------------------------------------------------------
% \archarrow{start}{end}{color} -- block arrow between stages
% ---------------------------------------------------------------
% #1 and #2 are full coordinates including parentheses, e.g. (1,2)
\newcommand{\archarrow}[3]{%
  \draw[-{Latex[length=2.4mm,width=2mm]},line width=2.2pt,#3] #1 -- #2;
}
 after \usepackage{tikz,ifthen} and
% \usetikzlibrary{arrows.meta,calc,positioning}. All macro names are
% prefixed `arch` to avoid clashing with pgf/tikz internals.

% ---- palette (approximating the PowerPoint drafts) ----
\definecolor{archTargetGreen}{HTML}{9BC7AC}
\definecolor{archContextPink}{HTML}{E8B8B0}
\definecolor{archAttnBlue}{HTML}{BBD6EC}
\definecolor{archEncGray}{HTML}{E4E4E4}
\definecolor{archArrowPale}{HTML}{CFE6D6}
\definecolor{archCtBg}{HTML}{2B2B2B}
\definecolor{archCtBlob}{HTML}{8C8C8C}
\definecolor{archCtxYellow}{HTML}{E8D24C}
% darker, more saturated siblings of archTargetGreen/archContextPink --
% pastel fills read poorly as 1pt strokes, and mask/label content needs to
% read as "which volume owns this" at a glance, so masks/labels and any
% line (arrow, brace) touching that volume's row use these instead of the
% pastel role color or the old flat archMaskGold.
\definecolor{archTargetAccent}{HTML}{2F6B4A}  % target: darker green
\definecolor{archContextAccent}{HTML}{A8453E} % context: darker red

% fake-isometric depth skew shared by every cube
\newcommand{\archdx}{0.30}
\newcommand{\archdy}{0.20}

% ---- shared text-role macros ----
% Every arch_*.tex figure is compiled standalone at the SAME base size
% (documentclass[tikz,border=4pt,12pt]{standalone}, matching the thesis
% body's 12pt), so \normalsize etc. mean the same absolute point size in
% every figure. Using these three named roles instead of ad hoc \Large /
% \bfseries\normalsize / \scriptsize per figure keeps "a box title" (etc.)
% the same size everywhere it appears, regardless of which figure it's in
% or how that figure's canvas happens to be scaled by \includegraphics.
%   \archtitle    -- figure/method/panel titles, the "xL" repeat marker
%   \archboxlabel -- sub-block header labels (Fuse, Encoder, attention, ...)
%   \archnote     -- small secondary annotations (placeholder captions)
% Plain, unlabeled node text (cube/token labels, row/lane labels) stays at
% the bare \normalsize default -- no macro needed for that role.
%
% DECLARATIONS, not argument-taking macros -- deliberately, so they can
% precede a raw `\\` inside a tikz node's text (as `{\archboxlabel line
% one\\line two}`) without hiding that `\\` one macro-expansion level deep.
% A `\\` buried inside a wrapped {\archboxlabel{...}} argument breaks
% TikZ's own align=center line-splitting (it scans for `\\` before the
% argument expands) -- use these the same way `\bfseries`/`\Large` are
% already used directly inline elsewhere in these figures.
\newcommand{\archtitle}{\Large\bfseries}
\newcommand{\archboxlabel}{\bfseries\normalsize}
\newcommand{\archnote}{\scriptsize}

% ---------------------------------------------------------------
% \archvolcube{x}{y}{w}{h}{topcolor}{content}{label}
%   content: ct | ctmask | predmask | <flat color name>
% ---------------------------------------------------------------
\newcommand{\archvolcube}[7]{%
  \begin{scope}
    \pgfmathsetmacro{\cx}{#1}\pgfmathsetmacro{\cy}{#2}%
    \pgfmathsetmacro{\w}{#3}\pgfmathsetmacro{\h}{#4}%
    \fill[#5] (\cx,\cy+\h) -- ++(\archdx,\archdy) -- ++(\w,0) -- ++(-\archdx,-\archdy) -- cycle;
    \fill[#5!60!black] (\cx+\w,\cy) -- ++(\archdx,\archdy) -- ++(0,\h) -- ++(-\archdx,-\archdy) -- cycle;
    \begin{scope}
      \clip (\cx,\cy) rectangle ++(\w,\h);
      \ifthenelse{\equal{#6}{ct}}{%
        \fill[archCtBg] (\cx,\cy) rectangle ++(\w,\h);
        \fill[archCtBlob] plot[smooth cycle] coordinates
          {(\cx+.25*\w,\cy+.55*\h) (\cx+.45*\w,\cy+.75*\h) (\cx+.7*\w,\cy+.6*\h)
           (\cx+.6*\w,\cy+.3*\h) (\cx+.35*\w,\cy+.25*\h)};
      }{\ifthenelse{\equal{#6}{ctmask}}{%
        \fill[archCtBg] (\cx,\cy) rectangle ++(\w,\h);
        \fill[archCtBlob] plot[smooth cycle] coordinates
          {(\cx+.25*\w,\cy+.55*\h) (\cx+.45*\w,\cy+.75*\h) (\cx+.7*\w,\cy+.6*\h)
           (\cx+.6*\w,\cy+.3*\h) (\cx+.35*\w,\cy+.25*\h)};
        \fill[archCtxYellow] plot[smooth cycle] coordinates
          {(\cx+.42*\w,\cy+.42*\h) (\cx+.52*\w,\cy+.5*\h) (\cx+.5*\w,\cy+.35*\h) (\cx+.44*\w,\cy+.3*\h)};
      }{\ifthenelse{\equal{#6}{predmask}}{%
        \fill[black] (\cx,\cy) rectangle ++(\w,\h);
        \fill[white] plot[smooth cycle] coordinates
          {(\cx+.4*\w,\cy+.65*\h) (\cx+.6*\w,\cy+.7*\h) (\cx+.62*\w,\cy+.45*\h)
           (\cx+.48*\w,\cy+.3*\h) (\cx+.38*\w,\cy+.4*\h)};
      }{%
        \fill[#6] (\cx,\cy) rectangle ++(\w,\h);
      }}}%
    \end{scope}
    \draw[black,thin] (\cx,\cy) rectangle ++(\w,\h);
    \draw[black,thin] (\cx,\cy+\h) -- ++(\archdx,\archdy) -- ++(\w,0) -- ++(-\archdx,-\archdy) -- cycle;
    \draw[black,thin] (\cx+\w,\cy) -- ++(\archdx,\archdy) -- ++(0,\h) -- ++(-\archdx,-\archdy) -- cycle;
    \node[anchor=north,align=center] at (\cx+\w/2,\cy-0.1) {#7};
  \end{scope}
}

% ---------------------------------------------------------------
% \archfeatcube{x}{y}{w}{h}{color}{label} -- flat-colored feature volume
% ---------------------------------------------------------------
\newcommand{\archfeatcube}[6]{%
  \begin{scope}
    \pgfmathsetmacro{\cx}{#1}\pgfmathsetmacro{\cy}{#2}%
    \pgfmathsetmacro{\w}{#3}\pgfmathsetmacro{\h}{#4}%
    \fill[#5] (\cx,\cy) rectangle ++(\w,\h);
    \fill[#5!80!white] (\cx,\cy+\h) -- ++(\archdx,\archdy) -- ++(\w,0) -- ++(-\archdx,-\archdy) -- cycle;
    \fill[#5!60!black] (\cx+\w,\cy) -- ++(\archdx,\archdy) -- ++(0,\h) -- ++(-\archdx,-\archdy) -- cycle;
    \draw[black,thin] (\cx,\cy) rectangle ++(\w,\h);
    \draw[black,thin] (\cx,\cy+\h) -- ++(\archdx,\archdy) -- ++(\w,0) -- ++(-\archdx,-\archdy) -- cycle;
    \draw[black,thin] (\cx+\w,\cy) -- ++(\archdx,\archdy) -- ++(0,\h) -- ++(-\archdx,-\archdy) -- cycle;
    \node[align=center] at (\cx+\w/2,\cy+\h/2) {#6};
  \end{scope}
}

% ---------------------------------------------------------------
% \archfunnel{x}{y}{h}{wwide}{wnarrow}{dir}{color}{label}
%   dir=1  : wide on left, narrow on right  (encoder)
%   dir=-1 : narrow on left, wide on right  (decoder)
% ---------------------------------------------------------------
\newcommand{\archfunnel}[8]{%
  \begin{scope}
    \pgfmathsetmacro{\cx}{#1}\pgfmathsetmacro{\cy}{#2}\pgfmathsetmacro{\h}{#3}%
    \pgfmathsetmacro{\ww}{#4}\pgfmathsetmacro{\wn}{#5}%
    \pgfmathsetmacro{\inset}{(\ww-\wn)/2}%
    \ifnum#6=1
      \fill[#7] (\cx,\cy) -- (\cx+\ww,\cy+\inset) -- (\cx+\ww,\cy+\h-\inset) -- (\cx,\cy+\h) -- cycle;
      \draw[black,thin] (\cx,\cy) -- (\cx+\ww,\cy+\inset) -- (\cx+\ww,\cy+\h-\inset) -- (\cx,\cy+\h) -- cycle;
    \else
      \fill[#7] (\cx,\cy+\inset) -- (\cx+\ww,\cy) -- (\cx+\ww,\cy+\h) -- (\cx,\cy+\h-\inset) -- cycle;
      \draw[black,thin] (\cx,\cy+\inset) -- (\cx+\ww,\cy) -- (\cx+\ww,\cy+\h) -- (\cx,\cy+\h-\inset) -- cycle;
    \fi
    \node[align=center] at (\cx+\ww/2,\cy+\h/2) {#8};
  \end{scope}
}

% ---------------------------------------------------------------
% \archbowtie{x}{y}{w}{h}{pinch}{color}{label} -- Medverse-style
% single-row "dual U-Net" hourglass (funnel-in then funnel-out).
% ---------------------------------------------------------------
\newcommand{\archbowtie}[7]{%
  \begin{scope}
    \pgfmathsetmacro{\cx}{#1}\pgfmathsetmacro{\cy}{#2}%
    \pgfmathsetmacro{\w}{#3}\pgfmathsetmacro{\h}{#4}\pgfmathsetmacro{\p}{#5}%
    \pgfmathsetmacro{\inset}{(\h-\p)/2}%
    \fill[#6] (\cx,\cy) -- (\cx+\w/2,\cy+\inset) -- (\cx+\w,\cy) -- (\cx+\w,\cy+\h)
      -- (\cx+\w/2,\cy+\h-\inset) -- (\cx,\cy+\h) -- cycle;
    \draw[black,thin] (\cx,\cy) -- (\cx+\w/2,\cy+\inset) -- (\cx+\w,\cy) -- (\cx+\w,\cy+\h)
      -- (\cx+\w/2,\cy+\h-\inset) -- (\cx,\cy+\h) -- cycle;
    \node[align=center] at (\cx+\w/2,\cy+\h/2) {#7};
  \end{scope}
}

% ---------------------------------------------------------------
% \archtokengrid{x}{y}{rows}{cols}{cell}{color} -- grid of patch tokens
% ---------------------------------------------------------------
\newcommand{\archtokengrid}[6]{%
  \begin{scope}
    \pgfmathsetmacro{\cx}{#1}\pgfmathsetmacro{\cy}{#2}\pgfmathsetmacro{\cell}{#5}%
    \pgfmathtruncatemacro{\rmax}{#3-1}%
    \pgfmathtruncatemacro{\cmax}{#4-1}%
    \foreach \r in {0,...,\rmax}{%
      \foreach \c in {0,...,\cmax}{%
        \fill[#6] ({\cx+\c*\cell},{\cy+\r*\cell}) rectangle ++(\cell*0.82,\cell*0.82);
        \draw[black,thin] ({\cx+\c*\cell},{\cy+\r*\cell}) rectangle ++(\cell*0.82,\cell*0.82);
      }%
    }%
  \end{scope}
}

% ---------------------------------------------------------------
% \archarrow{start}{end}{color} -- block arrow between stages
% ---------------------------------------------------------------
% #1 and #2 are full coordinates including parentheses, e.g. (1,2)
\newcommand{\archarrow}[3]{%
  \draw[-{Latex[length=2.4mm,width=2mm]},line width=2.2pt,#3] #1 -- #2;
}


% Define a vibrant, brighter blue for all attention arrows
\definecolor{archBrightBlue}{HTML}{3B82F6}

% ---------------------------------------------------------------------------
% Three fusion strategies, stacked so the pipelines read left-to-right
% within each panel and top-to-bottom across panels: Medverse
% (context-fused, cross-attn -- see arch_medverse_staircase.tex for this
% panel's standalone source and arch_medverse_attn_detail.tex for the
% per-block wiring its second box collapses), Baseline (early fusion,
% this thesis's single-axis control), Ours (dual attention). Style follows TabPFN's
% Fig. 1 (Hollmann et al., Nature 2024): small token grids connected by
% simple curved double-arrows rather than full cell/token-block
% renderings. The "special test row" device (edges touching it get a
% distinct color) maps onto our target row, which really is
% architecturally special in all three methods: it cross-attends/fuses
% read-only into context, and its label-stream input is either a
% predicted prior (Baseline/Ours) or entirely absent (Medverse) -- never
% real ground truth like context's, which is why target's label tokens
% are gray rather than a role color throughout. Each image/label volume
% is patchified into N spatial tokens; each cell below shows 3 of them as
% a representative sample (not literally N=3).
%
% Medverse (top): 1x Fuse, context only -- target's image tokens bypass
% it entirely, since Medverse never gives target a label channel at all,
% not even the zero-mask placeholder Baseline/Ours use -- followed by Lx
% Per-img conv block + cross-attn, where each row's own CNN block (own
% weights, no interaction) feeds directly into the token column that the
% cross-attention fan couples; same xL convention as Baseline, wrapping
% only the repeated box. Note L means something different here: U-Net
% stages, not transformer layers, and the real attention is asymmetric
% between the narrowing/widening halves of Medverse's two U-Nets --
% simplified here to one generic coupling per repeat.
% Baseline (middle): 1x Fuse (image/label tokens merged once per cell by
% one shared operator -- concat / sum / MLP / pixel-shuffle, drawn as a
% single ⊕) followed by Lx Patch axis -- fusion happens once, patch-axis
% attention repeats for L layers, so only the Patch-axis box carries the
% "xL" mark.
% Ours (bottom): the WHOLE pair (image/label axis -> patch axis) repeats
% as one layer, Lx -- so both boxes sit inside one outer dashed wrapper
% that carries the single "xL" mark for the pair.
%
% All boxes across all three panels share the SAME height, with each
% panel's content centered inside its box -- even though the baseline's
% single-stream patch-axis panel only fills a fraction of its box on the
% fan side, and Medverse's Fuse box leaves target's fuse arrow unused.
% That empty space is the point: the box size communicates "this step",
% not "how much content it happens to have". Medverse's second box is
% wider (\boxwB) than the other panels' boxes (\boxw), since it holds two
% sub-columns instead of one.
% ---------------------------------------------------------------------------

% \tpsq{x}{y}{fill}{ring} -- one token, a filled square with a role ring.
\newcommand{\tpsq}[4]{
  \fill[#3] (#1-\tsz/2,#2-\tsz/2) rectangle ++(\tsz,\tsz);
  \draw[#4,line width=1.2pt] (#1-\tsz/2,#2-\tsz/2) rectangle ++(\tsz,\tsz);
}
% \tpcluster{x}{y}{fill}{ring} -- 3 stacked tokens standing in for N.
\newcommand{\tpcluster}[4]{
  \foreach \doff in {-\toff,0,\toff} {
    \tpsq{#1}{#2+\doff}{#3}{#4}
  }
}
% \tpsqsplit{x}{y}{fillImg}{fillLbl}{ring} -- one token, diagonally split
\newcommand{\tpsqsplit}[5]{
  \begin{scope}
    \clip (#1-\tsz/2,#2-\tsz/2) rectangle ++(\tsz,\tsz);
    \fill[#3] (#1-\tsz/2,#2-\tsz/2) -- (#1+\tsz/2,#2-\tsz/2) -- (#1+\tsz/2,#2+\tsz/2) -- cycle;
    \fill[#4] (#1-\tsz/2,#2-\tsz/2) -- (#1-\tsz/2,#2+\tsz/2) -- (#1+\tsz/2,#2+\tsz/2) -- cycle;
  \end{scope}
  \draw[#5,line width=1.2pt] (#1-\tsz/2,#2-\tsz/2) rectangle ++(\tsz,\tsz);
}
% \tpclustersplit{x}{y}{fillImg}{fillLbl}{ring} -- 3 stacked split tokens.
\newcommand{\tpclustersplit}[5]{
  \foreach \doff in {-\toff,0,\toff} {
    \tpsqsplit{#1}{#2+\doff}{#3}{#4}{#5}
  }
}
% \dblarrow{p1}{p2}{color} -- curved double-headed attention edge.
\newcommand{\dblarrow}[3]{
  \draw[{Latex[length=3.0mm,width=2.5mm]}-{Latex[length=3.0mm,width=2.5mm]},
        line width=1.5pt,#3] #1 to[bend left=12] #2;
}
% \fusearrow{p1}{p2}{color} -- straight single-headed merge edge.
\newcommand{\fusearrow}[3]{
  \draw[-{Latex[length=3.0mm,width=2.5mm]},line width=1.5pt,#3] #1 -- #2;
}
% \cnnblock{x}{y}{fill}{draw} -- one row's own-weights conv block (Medverse
% panel only), drawn as a plain labeled "CNN" rectangle rather than a
% token cluster.
\newcommand{\cnnblock}[4]{
  \draw[#4,fill=#3,line width=1.2pt,rounded corners=4pt]
    (#1-\cnnw/2,#2-\cnnh/2) rectangle ++(\cnnw,\cnnh);
  \node[align=center] at (#1,#2) {\archboxlabel CNN};
}

% \tabpfnattn{x} -- Fully connected, right-side routed attention with cascading transparency
\newcommand{\tabpfnattn}[1]{
  
  % Base properties for all attention arrows using the new brighter blue
  \def\attnstyle{[{Latex[length=3.0mm,width=2.5mm]}-{Latex[length=3.0mm,width=2.5mm]}, line width=1.5pt, archBrightBlue]}
  \def\ax{#1+\tsz/2+0.05}

  % Distance 1: Adjacent neighbors (Opacity 1.0)
  \draw\attnstyle (\ax, \yTgt+\toff) to[bend left=30] (\ax, \yTgt);
  \draw\attnstyle (\ax, \yTgt)       to[bend left=30] (\ax, \yTgt-\toff);
  \draw\attnstyle (\ax, \yTgt-\toff) to[bend left=30] (\ax, \yCtx+\toff); % The gap bridge
  \draw\attnstyle (\ax, \yCtx+\toff) to[bend left=30] (\ax, \yCtx);
  \draw\attnstyle (\ax, \yCtx)       to[bend left=30] (\ax, \yCtx-\toff);

  % Distance 2: Skip 1 token (Opacity 0.6)
  \begin{scope}[opacity=0.6]
    \draw\attnstyle (\ax, \yTgt+\toff) to[bend left=45] (\ax, \yTgt-\toff);
    \draw\attnstyle (\ax, \yTgt)       to[bend left=45] (\ax, \yCtx+\toff);
    \draw\attnstyle (\ax, \yTgt-\toff) to[bend left=45] (\ax, \yCtx);
    \draw\attnstyle (\ax, \yCtx+\toff) to[bend left=45] (\ax, \yCtx-\toff);
  \end{scope}

  % Distance 3: Skip 2 tokens (Opacity 0.3)
  \begin{scope}[opacity=0.3]
    \draw\attnstyle (\ax, \yTgt+\toff) to[bend left=60] (\ax, \yCtx+\toff);
    \draw\attnstyle (\ax, \yTgt)       to[bend left=60] (\ax, \yCtx);
    \draw\attnstyle (\ax, \yTgt-\toff) to[bend left=60] (\ax, \yCtx-\toff);
  \end{scope}

  % Distance 4: Skip 3 tokens (Opacity 0.15)
  \begin{scope}[opacity=0.15]
    \draw\attnstyle (\ax, \yTgt+\toff) to[bend left=75] (\ax, \yCtx);
    \draw\attnstyle (\ax, \yTgt)       to[bend left=75] (\ax, \yCtx-\toff);
  \end{scope}

  % Distance 5: Skip 4 tokens - Full sequence wrap (Opacity 0.08)
  \begin{scope}[opacity=0.08]
    \draw\attnstyle (\ax, \yTgt+\toff) to[bend left=90] (\ax, \yCtx-\toff);
  \end{scope}
}

% ---- Expanded dimensions for HUGE tokens and continuous columns ----
\def\tsz{0.50}         
\def\toff{0.80}        
\def\rowpitch{2.40}    
\def\colpitch{2.4}     
\def\panelgap{2.6}     
\def\clusterhalf{1.10} 
\def\pairpitch{1.1}    

\def\xlab{-0.8}        
\def\yTgt{\rowpitch}   
\def\yCtx{0}

\def\boxw{6.0}
\def\xA{-0.30}
\def\ximgTop{1.0}      % Adjusted from 0.7 to keep Baseline columns centered
\def\boxpad{0.60}

% ---- Medverse panel only: its second box is wider than \boxw, since it
% holds two sub-columns (CNN block, token column) instead of one ----
\def\boxwB{7.6}
\def\cnnw{1.9}
\def\cnnh{1.9}

\begin{document}
\begin{tikzpicture}[x=1cm,y=1cm]

  % ---- pre-computed derived positions ----
  \pgfmathsetmacro{\xB}{\xA+\boxw+\panelgap}                  
  \pgfmathsetmacro{\boxAcenter}{\xA+\boxw/2}
  \pgfmathsetmacro{\boxBcenter}{\xB+\boxw/2}
  
  % Baseline wrapper coordinates
  \pgfmathsetmacro{\outerLeftBase}{\xB-0.50}
  \pgfmathsetmacro{\outerRightBase}{\xB+\boxw+0.50}
  
  % Stacked Layout: Ours shares X-coords with Baseline
  \pgfmathsetmacro{\xC}{\xA}                        
  \pgfmathsetmacro{\xD}{\xB}                   
  \pgfmathsetmacro{\boxCcenter}{\xC+\boxw/2}
  \pgfmathsetmacro{\boxDcenter}{\xD+\boxw/2}
  
  % Ours wrapper coordinates
  \pgfmathsetmacro{\outerLeft}{\xC-0.50}
  \pgfmathsetmacro{\outerRight}{\xD+\boxw+0.50}

  % Medverse layout: shares its left x-coord with Baseline/Ours, but its
  % second box is wider (\boxwB, not \boxw) to fit the CNN + token
  % sub-columns
  \pgfmathsetmacro{\xE}{\xA}
  \pgfmathsetmacro{\xF}{\xE+\boxw+\panelgap}
  \pgfmathsetmacro{\boxEcenter}{\xE+\boxw/2}
  \pgfmathsetmacro{\boxFcenter}{\xF+\boxwB/2}
  \pgfmathsetmacro{\xCnnCol}{\xF+1.6}
  \pgfmathsetmacro{\xTokCol}{\xF+\boxwB-1.6}

  % Medverse wrapper coordinates -- spans only the second box (Per-img
  % conv block + cross-attn), matching how Baseline's wrapper spans only
  % its Patch axis box: Fuse happens once, upstream, in both panels
  \pgfmathsetmacro{\outerLeftMV}{\xF-0.50}
  \pgfmathsetmacro{\outerRightMV}{\xF+\boxwB+0.50}

  % Shared vertical boundaries for outer wrappers (Local Y=0)
  \pgfmathsetmacro{\outerTop}{\yTgt+\clusterhalf+\boxpad+0.85}
  \pgfmathsetmacro{\outerBot}{\yCtx-\clusterhalf-\boxpad-0.40}
  
  \pgfmathsetmacro{\methodTitleY}{\outerTop+0.50}              
  
  % X placements for columns inside boxes
  \pgfmathsetmacro{\xlblTop}{\ximgTop+1.4}     % Reduced gap for Baseline img and lbl columns
  \pgfmathsetmacro{\xFusedTop}{\xlblTop+2.0}
  
  % Fuse operator dynamically placed exactly between img and lbl columns
  \pgfmathsetmacro{\xFuseOp}{(\ximgTop+\xlblTop)/2}
  \pgfmathsetmacro{\xFuseArrowStart}{\xFuseOp+0.35} % Adjusted for bare oplus node
  \pgfmathsetmacro{\yMid}{\rowpitch/2}
  
  \pgfmathsetmacro{\xFusedPatchTop}{\boxBcenter}
  \pgfmathsetmacro{\ximgBot}{\boxCcenter-\colpitch/2}
  \pgfmathsetmacro{\xlblBot}{\ximgBot+\colpitch}
  \pgfmathsetmacro{\ximgPatchBot}{\boxDcenter-\pairpitch/2}
  \pgfmathsetmacro{\xlblPatchBot}{\ximgPatchBot+\pairpitch}


  % =========================================================================
  % Medverse: 1x Fuse (context only), then Lx Per-img conv block +
  % cross-attn (TOP BLOCK). See arch_medverse_staircase.tex for the
  % standalone version this is lifted from, and arch_medverse_attn_detail.tex
  % for the literal per-block wiring the second box collapses.
  % =========================================================================
  \begin{scope}[yshift=0cm]

    % Row labels for Medverse
    \node[anchor=east,font=\normalsize] at (\xlab,\yTgt) {target ($N$)};
    \node[anchor=east,font=\normalsize] at (\xlab,\yCtx) {context ($N$)};

    % Title for Medverse -- named for its architecture family (dual,
    % separately-weighted CNNs), matching how block 2 below is named for
    % its own family rather than for either individual method
    \node[anchor=south] at ({\xE+(\xF+\boxwB-\xE)/2},{\methodTitleY})
      {\archtitle Dual ConvNet \normalsize\normalfont(Medverse)};

    % Medverse outer wrapper (Lx repetition marker) -- spans only the
    % second box, same convention as Baseline's wrapper: Fuse happens
    % once, upstream, in both panels
    \draw[archTargetGreen!60!black,fill=archTargetGreen!10,dashed,rounded corners=8pt,line width=1pt]
      ({\outerLeftMV},{\outerBot}) rectangle ({\outerRightMV},{\outerTop});

    % Inner boxes (filled white)
    \draw[black!30,fill=white,rounded corners=5pt,line width=0.8pt]
      (\xE,\yCtx-\clusterhalf-\boxpad) rectangle ({\xE+\boxw},\yTgt+\clusterhalf+\boxpad);
    \draw[black!30,fill=white,rounded corners=5pt,line width=0.8pt]
      ({\xF},\yCtx-\clusterhalf-\boxpad) rectangle ({\xF+\boxwB},\yTgt+\clusterhalf+\boxpad);

    \node[anchor=south] at ({\boxEcenter},\yTgt+\clusterhalf+\boxpad+0.20) {\archboxlabel Fuse (context only)};
    \node[anchor=south,align=center] at ({\boxFcenter},\yTgt+\clusterhalf+\boxpad+0.20)
      {\archboxlabel Per-img conv block + cross-attn};

    % "xL" placed in the top right corner of the wrapper
    \node[anchor=north east] at ({\outerRightMV-0.20},{\outerTop-0.20}) {\archtitle $\times L$};

    % Arch arrow connecting the two steps
    \archarrow{({\xE+\boxw+0.50},{\yMid})}{({\xF-0.50},{\yMid})}{black!30}

    \node[anchor=south,font=\normalsize] at (\ximgTop,\yTgt+\clusterhalf+0.25) {img};
    \node[anchor=south,font=\normalsize] at (\xlblTop,\yTgt+\clusterhalf+0.25) {lbl};
    \node[anchor=south,font=\normalsize] at (\xFusedTop,\yTgt+\clusterhalf+0.25) {fused};

    % Background columns for Fuse (Vertical grouping)
    \foreach \x in {\ximgTop, \xlblTop, \xFusedTop} {
      \fill[black!15,rounded corners=3pt]
        ({\x-\tsz/2-0.18},{\yCtx-\toff-\tsz/2-0.10}) rectangle
        ({\x+\tsz/2+0.18},{\yTgt+\toff+\tsz/2+0.10});
    }

    % target row: raw image tokens, unsplit fill, bypass the fuse. Unlike
    % the other two panels, target here has no label column at all -- not
    % even a zero-mask placeholder, Medverse's TargetEncoder never
    % receives one. The gray lbl token is kept anyway, purely for column
    % alignment with the other two panels; it is deliberately left with
    % no arrows in or out, since it plays no role in the fuse
    \tpcluster{\ximgTop}{\yTgt}{archTargetGreen}{black!70}
    \tpcluster{\xlblTop}{\yTgt}{black!8}{black!70}
    \tpcluster{\xFusedTop}{\yTgt}{archTargetGreen}{black!70}
    \fusearrow{({\ximgTop+\tsz/2+0.18},{\yTgt})}{({\xFusedTop-\tsz/2-0.18},{\yTgt})}{black!40}

    % context row: img + lbl -> (+) -> fused
    \tpcluster{\ximgTop}{\yCtx}{archContextPink}{black!70}
    \tpcluster{\xlblTop}{\yCtx}{archContextAccent}{black!70}
    \tpclustersplit{\xFusedTop}{\yCtx}{archContextPink}{archContextAccent}{black!70}
    \node[font=\huge] at ({\xFuseOp},{\yCtx}) {$\oplus$};
    \fusearrow{({\xFuseArrowStart},{\yCtx})}{({\xFusedTop-\tsz/2-0.18},{\yCtx})}{black!55}

    % Per-img conv block + cross-attn: CNN block (own weights, no
    % interaction) sits right next to the token column it feeds
    \cnnblock{\xCnnCol}{\yTgt}{archTargetGreen!35}{archTargetAccent}
    \cnnblock{\xCnnCol}{\yCtx}{archContextPink!35}{archContextAccent}
    \fusearrow{({\xCnnCol+\cnnw/2+0.15},{\yTgt})}{({\xTokCol-\tsz/2-0.18},{\yTgt})}{black!40}
    \fusearrow{({\xCnnCol+\cnnw/2+0.15},{\yCtx})}{({\xTokCol-\tsz/2-0.18},{\yCtx})}{black!40}

    % Background row groupings (token column)
    \foreach \y in {\yTgt,\yCtx} {
      \foreach \doff in {-\toff,0,\toff} {
        \fill[black!15,rounded corners=3pt]
          ({\xTokCol-\tsz/2-0.18},{\y+\doff-\tsz/2-0.10}) rectangle
          ({\xTokCol+\tsz/2+0.18},{\y+\doff+\tsz/2+0.10});
      }
    }

    % token column: the only place the two streams actually meet
    \tpcluster{\xTokCol}{\yTgt}{archTargetGreen}{black!70}
    \tpclustersplit{\xTokCol}{\yCtx}{archContextPink}{archContextAccent}{black!70}
    \tabpfnattn{\xTokCol}

  \end{scope}


  % =========================================================================
  % Horizontal Divider (between Medverse and Baseline). Spans the widest
  % panel's extent (Medverse's) so both dividers line up regardless of
  % which panel sits in which slot.
  % =========================================================================
  \pgfmathsetmacro{\divY}{\outerBot - 1.2}
  \draw[black!20,dashed,line width=1pt] ({\xA-1.0},{\divY}) -- ({\xF+\boxwB+1.0},{\divY});


  % =========================================================================
  % Block 2: "Shared encoder + cross-attn" -- Baseline and Ours are its two
  % sub-blocks, both built on the same encoder+attention family (only the
  % fusion axis differs between them), unlike block 1's dual, separately-
  % weighted CNNs. Baseline: 1x Fuse, then Lx Patch axis (MIDDLE BLOCK)
  % =========================================================================
  \begin{scope}[yshift=-10.3cm]

    % Row labels for Baseline
    \node[anchor=east,font=\normalsize] at (\xlab,\yTgt) {target ($N$)};
    \node[anchor=east,font=\normalsize] at (\xlab,\yCtx) {context ($N$)};

    % Group title, spanning both sub-blocks (Baseline + Ours below) --
    % sits above Baseline's own title, in the gap between the block 1/2
    % divider and Baseline's content
    \node[anchor=south] at ({\xA+(\xB+\boxw-\xA)/2},{\methodTitleY+0.80})
      {\archtitle Shared encoder + cross-attn};

    % Title for Baseline (centered across both inner boxes)
    \node[anchor=south] at ({\xA+(\xB+\boxw-\xA)/2},{\methodTitleY}) {\archboxlabel Baseline: early fusion};

    % Baseline outer wrapper (Lx repetition marker)
    \draw[archTargetGreen!60!black,fill=archTargetGreen!10,dashed,rounded corners=8pt,line width=1pt]
      ({\outerLeftBase},{\outerBot}) rectangle ({\outerRightBase},{\outerTop});

    % Inner boxes (filled white to sit cleanly above the green wrapper)
    \draw[black!30,fill=white,rounded corners=5pt,line width=0.8pt]
      (\xA,\yCtx-\clusterhalf-\boxpad) rectangle ({\xA+\boxw},\yTgt+\clusterhalf+\boxpad);
    \draw[black!30,fill=white,rounded corners=5pt,line width=0.8pt]
      ({\xB},\yCtx-\clusterhalf-\boxpad) rectangle ({\xB+\boxw},\yTgt+\clusterhalf+\boxpad);

    \node[anchor=south] at ({\boxAcenter},\yTgt+\clusterhalf+\boxpad+0.20) {\archboxlabel Fuse};
    \node[anchor=south] at ({\boxBcenter},\yTgt+\clusterhalf+\boxpad+0.20) {\archboxlabel Patch axis};

    % "xL" placed in the top right corner of the wrapper
    \node[anchor=north east] at ({\outerRightBase-0.20},{\outerTop-0.20}) {\archtitle $\times L$};

    % Arch arrow connecting the two steps
    \archarrow{({\xA+\boxw+0.50},{\yMid})}{({\xB-0.50},{\yMid})}{black!30}

    \node[anchor=south,font=\normalsize] at (\ximgTop,\yTgt+\clusterhalf+0.25) {img};
    \node[anchor=south,font=\normalsize] at (\xlblTop,\yTgt+\clusterhalf+0.25) {lbl};
    \node[anchor=south,font=\normalsize] at (\xFusedTop,\yTgt+\clusterhalf+0.25) {fused};
    \node[anchor=south,font=\normalsize] at ({\xFusedPatchTop},\yTgt+\clusterhalf+0.25) {fused};

    % Background columns for Fuse (Vertical grouping)
    \foreach \x in {\ximgTop, \xlblTop, \xFusedTop} {
      \fill[black!15,rounded corners=3pt]
        ({\x-\tsz/2-0.18},{\yCtx-\toff-\tsz/2-0.10}) rectangle
        ({\x+\tsz/2+0.18},{\yTgt+\toff+\tsz/2+0.10});
    }

    % Background rows for Patch axis (Horizontal grouping, single column)
    \foreach \y in {\yTgt,\yCtx} {
      \foreach \doff in {-\toff,0,\toff} {
        \fill[black!15,rounded corners=3pt]
          ({\xFusedPatchTop-\tsz/2-0.18},{\y+\doff-\tsz/2-0.10}) rectangle
          ({\xFusedPatchTop+\tsz/2+0.18},{\y+\doff+\tsz/2+0.10});
      }
    }

    % target's label column is gray, not a role color -- it is never real
    % ground truth (a zero-mask placeholder here; absent entirely in
    % Medverse), unlike context's, so it should not read as informative
    % content
    \foreach \y/\imgfill/\lblfill in {\yTgt/archTargetGreen/black!8,
                                       \yCtx/archContextPink/archContextAccent} {
      \tpcluster{\ximgTop}{\y}{\imgfill}{black!70}
      \tpcluster{\xlblTop}{\y}{\lblfill}{black!70}
      \tpclustersplit{\xFusedTop}{\y}{\imgfill}{\lblfill}{black!70}
      \tpclustersplit{\xFusedPatchTop}{\y}{\imgfill}{\lblfill}{black!70}
    }

    % Centered fuse operator without the circle
    \node[font=\huge] at ({\xFuseOp},{\yMid}) {$\oplus$};
    \fusearrow{({\xFuseArrowStart},{\yMid})}{({\xFusedTop-\tsz/2-0.18},{\yMid})}{black!55}

    % TabPFN-style attention routing for the Baseline
    \tabpfnattn{\xFusedPatchTop}

  \end{scope}


  % =========================================================================
  % Divider between the two sub-blocks (Baseline/Ours) -- 9.5cm below
  % Baseline's own yshift, same fixed sub-block pitch as before; only the
  % block1->block2 gap above grew to fit the new group title
  % =========================================================================
  \pgfmathsetmacro{\divYTwo}{\divY-10.3}
  \draw[black!20,dashed,line width=1pt] ({\xA-1.0},{\divYTwo}) -- ({\xF+\boxwB+1.0},{\divYTwo});


  % =========================================================================
  % Ours: Lx (image/label axis -> patch axis) -- block 2's second
  % sub-block (BOTTOM BLOCK)
  % =========================================================================
  \begin{scope}[yshift=-19.8cm]

    % Row labels for Ours
    \node[anchor=east,font=\normalsize] at (\xlab,\yTgt) {target ($N$)};
    \node[anchor=east,font=\normalsize] at (\xlab,\yCtx) {context ($N$)};

    % Title for Ours (sub-block title, same tier as Baseline's)
    \node[anchor=south] at ({\xC+(\xD+\boxw-\xC)/2},{\methodTitleY}) {\archboxlabel Ours: dual attn.};

    % Ours outer wrapper (Lx repetition marker)
    \draw[archTargetGreen!60!black,fill=archTargetGreen!10,dashed,rounded corners=8pt,line width=1pt]
      ({\outerLeft},{\outerBot}) rectangle ({\outerRight},{\outerTop});

    % Inner boxes (filled white)
    \draw[black!30,fill=white,rounded corners=5pt,line width=0.8pt]
      (\xC,\yCtx-\clusterhalf-\boxpad) rectangle ({\xC+\boxw},\yTgt+\clusterhalf+\boxpad);
    \draw[black!30,fill=white,rounded corners=5pt,line width=0.8pt]
      ({\xD},\yCtx-\clusterhalf-\boxpad) rectangle ({\xD+\boxw},\yTgt+\clusterhalf+\boxpad);

    \node[anchor=south] at ({\boxCcenter},\yTgt+\clusterhalf+\boxpad+0.20) {\archboxlabel Image/label axis};
    \node[anchor=south] at ({\boxDcenter},\yTgt+\clusterhalf+\boxpad+0.20) {\archboxlabel Patch axis};

    % "xL" placed in the top right corner of the wrapper
    \node[anchor=north east] at ({\outerRight-0.20},{\outerTop-0.20}) {\archtitle $\times L$};

    % Arch arrow connecting the two steps inside the wrapper
    \archarrow{({\xC+\boxw+0.50},{\yMid})}{({\xD-0.50},{\yMid})}{black!30}

    \node[anchor=south,font=\normalsize] at ({\ximgBot},\yTgt+\clusterhalf+0.25) {img};
    \node[anchor=south,font=\normalsize] at ({\xlblBot},\yTgt+\clusterhalf+0.25) {lbl};
    \node[anchor=south,font=\normalsize] at ({\ximgPatchBot},\yTgt+\clusterhalf+0.25) {img};
    \node[anchor=south,font=\normalsize] at ({\xlblPatchBot},\yTgt+\clusterhalf+0.25) {lbl};

    % Background columns for Image/label axis (Vertical grouping)
    \fill[black!15,rounded corners=3pt]
      ({\ximgBot-\tsz/2-0.18},{\yCtx-\toff-\tsz/2-0.10}) rectangle
      ({\ximgBot+\tsz/2+0.18},{\yTgt+\toff+\tsz/2+0.10});
    \fill[black!15,rounded corners=3pt]
      ({\xlblBot-\tsz/2-0.18},{\yCtx-\toff-\tsz/2-0.10}) rectangle
      ({\xlblBot+\tsz/2+0.18},{\yTgt+\toff+\tsz/2+0.10});

    % Background rows for Patch axis (Horizontal grouping)
    \foreach \y in {\yTgt,\yCtx} {
      \foreach \doff in {-\toff,0,\toff} {
        \fill[black!15,rounded corners=3pt]
          ({\ximgPatchBot-\tsz/2-0.18},{\y+\doff-\tsz/2-0.10}) rectangle
          ({\xlblPatchBot+\tsz/2+0.18},{\y+\doff+\tsz/2+0.10});
      }
    }

    % target's label column is gray, not a role color -- see the matching
    % note in the Baseline block above
    \foreach \y/\imgfill/\lblfill in {\yTgt/archTargetGreen/black!8,
                                       \yCtx/archContextPink/archContextAccent} {
      \tpcluster{\ximgBot}{\y}{\imgfill}{black!70}
      \tpcluster{\xlblBot}{\y}{\lblfill}{black!70}
      \tpcluster{\ximgPatchBot}{\y}{\imgfill}{black!70}
      \tpcluster{\xlblPatchBot}{\y}{\lblfill}{black!70}
    }

    \foreach \doff in {-\toff,0,\toff} {
      % Using archBrightBlue for horizontal image/label attention
      \dblarrow{({\ximgBot+\tsz/2+0.10},\yCtx+\doff)}{({\xlblBot-\tsz/2-0.10},\yCtx+\doff)}{archBrightBlue}
      \dblarrow{({\ximgBot+\tsz/2+0.10},\yTgt+\doff)}{({\xlblBot-\tsz/2-0.10},\yTgt+\doff)}{archBrightBlue}
    }

    % TabPFN-style attention routing for Ours (Label column only)
    \tabpfnattn{\xlblPatchBot}

  \end{scope}

\end{tikzpicture}
\end{document}